\chapter{Formy normalne macierzy zespolonych}

W niniejszym rozdziale zajmiemy się formami normalnymi macierzy zespolonych. Głównym celem jest przedstawienie macierzy za pomocą prostszych postaci, które pozwalają lepiej zrozumieć ich strukturę oraz własności.

Przed sformułowaniem twierdzeń dotyczących postaci normalnych macierzy zespolonych musimy omówić funkcję od macierzy oraz wielomian interpolacyjny Lagrange'a--Sylwestra, a także udowodnić kilka twierdzeń dotyczących rozkładu biegunowego. Pojęcia te będą nam potrzebne w dalszych rozważaniach i dowodach.

\section{Funkcja od macierzy i wielomian Lagrange'a--Sylwestra}
W tym podrozdziale spróbujemy zdefiniować, czym jest funkcja od macierzy $f(A)$, gdzie $A$ jest macierzą kwadratową wymiaru $n \times n$. Dalej w podrozdziale, wszędzie gdzie piszemy $A$, zakładamy, że jest to macierz kwadratowa wymiaru $n \times n$, jeśli nie powiedziano inaczej.

Najprostszym przypadkiem funkcji od macierzy jest wielomian od macierzy:
\begin{definition}[Wielomian od macierzy]
    Niech $A$ będzie macierzą kwadratową wymiarów $n \times n$, a $f(\lambda)$ dowolnym wielomianem $l$-tego stopnia zmiennej $\lambda \in \C$:
    $$f(\lambda) = a_0\lambda^{l} + a_1\lambda^{l - 1} + \ldots + a_{l - 1}\lambda + a_l, \quad l \in \N, \ a_i \in \C$$
    Wtedy $f(A)$ jest \textit{wielomianem od macierzy}:
    $$f(A) = a_0A^{l} + a_1A^{l - 1} + \ldots + a_{l - 1}A + a_lI, \quad l \in \N, \ a_i \in \C$$
\end{definition}
W dalszym ciągu spróbujemy sformułować postać ogólną funkcji od macierzy $f(A)$, ale najpierw podamy kilka definicji oraz twierdzeń. 
\begin{definition}\label{def:matrix_spectrum}
    Spektrum macierzy $A$ nazywamy zbiór wszystkich wartości własnych danej macierzy i oznaczamy przez $\Lambda_{A}$.
\end{definition}

\begin{definition}[Wielomian zerujący]
     Wielomianem zerującym macierzy $A$ będziemy nazywać taki wielomian $g(\lambda)$, dla którego zachodzi:
     $$g(A) = 0$$
\end{definition}
Na mocy Twierdzenia Cayleya--Hamiltona \cite[Rozdział IV, \S 4, str. 83]{gantmacher2000} wielomian charakterystyczny macierzy $A$ jest wielomianem zerującym.
\begin{definition}[Wielomian minimalny]
    Wielomian zerujący o najniższym stopniu i współczynniku wiodącym równym 1 będziemy nazywać \textit{wielomianem minimalnym}.
\end{definition}
\begin{theorem}[O wielomianie zerującym]\label{theorem:annihilating_polynomial}
    Dowolny wielomian zerujący jest podzielny przez wielomian minimalny.
\end{theorem}
\begin{proof}
Niech dana będzie macierz kwadratowa $A$ oraz niech $f(\lambda)$ będzie dowolnym wielomianem zerującym macierzy $A$ (tzn. $f(A) = 0$). Podzielmy $f(\lambda)$ przez wielomian minimalny $\psi(\lambda)$:
$$ f(\lambda) = \psi(\lambda)q(\lambda) + r(\lambda) $$
Z tego wynika:
$$ f(A) = \psi(A)q(A) + r(A) $$
Jak wiemy, $\psi(A) = 0$ oraz $f(A) = 0$:
$$ f(A) = r(A) \implies r(A) = 0 $$
Stopień $r(\lambda)$ jest mniejszy od stopnia wielomianu minimalnego, więc $r(\lambda) \equiv 0$, ponieważ w przeciwnym przypadku wielomian $r(\lambda)$ byłby wielomianem zerującym o mniejszym stopniu niż wielomian minimalny, co daje sprzeczność.
Zatem reszta jest zawsze zerowa, $r(\lambda) \equiv 0$:
$$f(\lambda) = \psi(\lambda)q(\lambda) \implies \psi(\lambda) \mid f(\lambda)$$
\end{proof}


Odnotujmy ważną własność wiążącą wielomian charakterystyczny oraz minimalny:
Niech $\varphi(\lambda)$ wielomian charakterystyczny macierzy $A$ oraz $\psi(\lambda)$ jest wielomianem minimalnym. 

Z Twierdzenia Cayleya--Hamiltona oraz Twierdzenia~\ref{theorem:annihilating_polynomial} (\textit{O wielomianie zerującym}) wynika, że wielomian charakterystyczny jest podzielny przez wielomian minimalny, więc ma co najmniej tyle samo różnych pierwiastków co wielomian minimalny:
$$ \varphi(\lambda) = \psi(\lambda)c(\lambda), $$
gdzie $c(\lambda)$ pewien wielomian od $\lambda$. A wielomian minimalny $\psi(\lambda)$ zeruje się dla każdego pierwiastka $\varphi(\lambda)$, tj. dla każdej wartości własnej. Weźmy dowolny wektor własny i wartość własną:
$$ Ax = \lambda x \implies \psi(A)x = \psi(\lambda)x \implies 0 = \psi(\lambda)x$$
% TODO: dowód że możemy pomnożyć na \psi
Ponieważ $x \ne 0$:
$$0 = \psi(\lambda)x \implies \psi(\lambda) = 0$$
Jak widzimy, rzeczywiście, wielomian minimalny zeruje się dla każdego pierwiastka wielomianu charakterystycznego, a więc różnica między wielomianem minimalnym i charakterystycznym jest tylko w krotności pierwiastków. Podsumujmy:

Wielomian minimalny $\psi(\lambda)$ ma te same różne pierwiastki co wielomian charakterystyczny, tj. wszystkie wartości własne macierzy $A$.
Jeśli:
\begin{equation}
   (\text{w. charakter.}) \quad 
   \varphi(\lambda) = (\lambda - \lambda_1)^{n_1}(\lambda - \lambda_2)^{n_2}\ldots(\lambda - \lambda_s)^{n_s}, \quad \sum_{k=1}^{s}{n_k} = n 
\end{equation}
to:
\begin{equation}\label{eq:minimal_poly}
    (\text{w. minimalny}) \quad 
    \psi(\lambda) = (\lambda - \lambda_1)^{m_1}(\lambda - \lambda_2)^{m_2}\ldots(\lambda - \lambda_s)^{m_s}, \quad \sum_{k=1}^{s}{m_k} = m
\end{equation}
gdzie $\lambda_1, \ldots, \lambda_s$ wszystkie różne wartości własne $A$, liczba $n_k$ reprezentuje krotność wartości własnej oraz $0 < m_k \le n_k \implies m \le n$.

\begin{definition}\label{def:function_on_spectrum}
    Niech $A$ będzie macierzą kwadratową, która ma następujące wartości własne:
    $$ \lambda^{(m_1)}_1, \ \lambda^{(m_2)}_2, \ \ldots, \ \lambda^{(m_s)}_s \quad (s = 1, 2, \ldots, n) $$
    gdzie $m_1, \ldots, m_s$ są wykładnikami wielomiana minimalnego \eqref{eq:minimal_poly}.

    Następujące liczby:
    \begin{equation}\label{values_on_spectrum}
        f(\lambda_{k}), \ f'(\lambda_{k}), \ \ldots, \ f^{(m_k - 1)}(\lambda_{k}) \quad (k = 1, 2, \ldots, s)
    \end{equation}
    będziemy nazywać wartościami na \textit{spektrum macierzy} $A$, a zbiór tych wartości będziemy oznaczać jako $f(\Lambda_{A})$. Jeśli dla funkcji $f(\lambda)$ istnieją wartości \eqref{values_on_spectrum}, to będziemy mówić, że funkcja jest zdefiniowana na spektrum macierzy.
\end{definition}

\begin{definition}[Wielomian Lagrange'a--Sylwestra]
    Niech $A$ będzie macierzą kwadratową, która ma następujący wielomian minimalny:
    $$ \psi(\lambda) = (\lambda - \lambda_1)^{m_1}(\lambda - \lambda_2)^{m_2}\ldots(\lambda - \lambda_s)^{m_s}, \quad \sum_{k=1}^{s}{m_k} = m$$
    gdzie $\lambda_1, \ldots, \lambda_s$ są wartościami własnymi $A$, liczba $m_k$ reprezentuje krotność czynnika w wielomianie minimalnym oraz $m \le n$.

    Niech dana będzie funkcja $f(\lambda)$ zdefiniowana na spektrum macierzy $A$, wtedy wielomianem Lagrange'a--Sylwestra będziemy nazywać jedyny wielomian interpolacyjny $r(\lambda)$ stopnia $< m$, który spełnia warunki:
    $$ f(\lambda_k) = r(\lambda_k), \ f'(\lambda_k) = r'(\lambda_k), \ \ldots, \ f^{(m_k - 1)}(\lambda_k) = r^{(m_k - 1)}(\lambda_k), \quad (k = 1, \ldots, s) $$
    Czyli $f(\lambda)$ i $r(\lambda)$ mają te same wartości oraz te same pochodne do rzędu $(m_k - 1)$ w każdym punkcie $\lambda_k$.
\end{definition}
\begin{theorem}\label{theorem:exist_L_S_poly}
Wielomian Lagrange'a--Sylwestra zawsze istnieje i jest jednoznacznie określony dla $f(\lambda)$.
\end{theorem}
\begin{proof}
Dowód pomijamy. Pełne uzasadnienie istnienia i jednoznaczności wielomianu Lagrange'a--Sylwestra można znaleźć w \cite[Rozdział V, \S 2, str. 103]{gantmacher2000}.
\end{proof}
% TODO: Pokazać, jak wygląda
\begin{theorem}
    Niech $A$ będzie macierzą kwadratową. Wszystkie wielomiany $g(\lambda)$, które mają te same wartości na spektrum macierzy $A$, mają tę samą wartość $g(A)$, tzn. jeśli dla dowolnych wielomianów $g_1(\lambda), \ldots, g_p(\lambda)$ zachodzi $g_1(\Lambda_{A}) = g_2(\Lambda_{A}) = \ldots = g_p(\Lambda_{A})$, to $g_1(A) = g_2(A) = \ldots = g_p(A)$.
\end{theorem}
\begin{proof}
    Oznaczmy przez $\psi(\lambda)$ wielomian minimalny macierzy $A$:
    \begin{equation}\label{eq:local_min_poly}
        \psi(\lambda) = (\lambda - \lambda_1)^{m_1}(\lambda - \lambda_2)^{m_2}\ldots(\lambda - \lambda_s)^{m_s}, \quad \sum_{k=1}^{s}{m_k} = m
    \end{equation}
    gdzie $\lambda_1, \ldots, \lambda_s$ są wszystkimi różnymi wartościami własnymi macierzy $A$, a $m$ jest stopniem wielomianu $\psi(\lambda)$ ($m \le n$).
    Niech $g(\lambda)$ i $h(\lambda)$ będą dowolnymi wielomianami takimi, że:
    \begin{equation}\label{eq:local_first_step}
        g(A) = h(A)
    \end{equation}
    Wtedy różnica $d(\lambda) = g(\lambda) - h(\lambda)$ jest wielomianem zerującym dla macierzy $A$ oraz na mocy Twierdzenia~\ref{theorem:annihilating_polynomial} (\textit{O wielomianie zerującym}) jest podzielny przez wielomian minimalny $\psi(\lambda)$. Stąd wynika, że wielomian $d(\lambda)$ ma co najmniej te same miejsca zerowe co $\psi(\lambda)$, w tym przypadku są to wartości własne. Możemy to łatwo sprawdzić, podstawiając dowolną wartość własną $\lambda_k$ do $d(\lambda)$, zawsze jakiś składnik z \eqref{eq:local_min_poly} wyzeruje się: $(\lambda - \lambda_k)^{m_k} = 0$.
    
    Tak samo jest w przypadku poszczególnych pochodnych $d'(\lambda), \ \ldots, \ d^{(m_k - 1)}(\lambda)$ do $m_k - 1$ rzędu, gdzie $k = 1, \ldots, s$, ponieważ każdy składnik $(\lambda - \lambda_k)^{m_k}$ możemy zróżniczkować $m_k - 1$ razy, nie sprowadzając tego składnika do stałej, więc $d'(\lambda), \ \ldots, \ d^{(m_k - 1)}(\lambda)$ będą zerować się dla każdej wartości własnej. Zapiszmy to jako:
    \begin{equation}\label{eq:local_second_step}
        d(\lambda_k) = 0, \ d'(\lambda_k) = 0, \ \ldots, \ d^{(m_k - 1)}(\lambda_k) = 0, \quad (k = 1, \ldots, s)
    \end{equation}
    Ponieważ różniczkowanie jest liniowe, możemy to przepisać jako:
    \begin{equation}\label{eq:local_third_step}
        g(\lambda_k) = h(\lambda_k), \ g'(\lambda_k) = h'(\lambda_k), \ \ldots, \ g^{(m_k - 1)}(\lambda_k) = h^{(m_k - 1)}(\lambda_k), \quad (k = 1, \ldots, s)
    \end{equation}
    Jak widzimy, wielomiany $g(\lambda)$ i $h(\lambda)$ przyjmują jednakowe wartości na spektrum macierzy $A$. 
    
    Zauważmy, że rozumowanie to działa w obie strony: z \eqref{eq:local_third_step} wynika \eqref{eq:local_second_step}, co jest oczywiste. Z \eqref{eq:local_second_step} możemy wyprowadzić \eqref{eq:local_first_step}: 
    Wielomian $d(\lambda)$ ma w każdym z punktów $\lambda_k$ zero o krotności co najmniej $m_k$ (tak jak wielomian minimalny $\psi(\lambda)$), więc $\psi(\lambda)$ dzieli $d(\lambda)$. Z tego wynika, że $d(\lambda)$ jest wielomianem zerującym dla $A$, wtedy:
$$d(A) = 0 \implies g(A) - h(A) = 0 \implies g(A) = h(A).$$

    Zatem wielomiany, które przyjmują jednakowe wartości na spektrum macierzy, mają jednakową wartość macierzową $g(A)$.
\end{proof}
Zdefiniujmy funkcję $f(A)$ analogicznie do powyższego twierdzenia, tzn. wszystkie funkcje $f(\lambda)$, które mają te same wartości na spektrum macierzy $A$, mają tę samą wartość $f(A)$.

Wtedy jest oczywiste, że dla definicji $f(A)$ możemy znaleźć wielomian $g(\lambda)$, który miałby te same wartości na spektrum macierzy $A$ co $f(\lambda)$, i przyjąć, że:
$$f(A) := g(A)$$
Na mocy Twierdzenia~\ref{theorem:exist_L_S_poly} wiemy, że wśród wszystkich tych wielomianów istnieje jeden i tylko jeden wielomian stopnia mniejszego od $m$, który jest wielomianem Lagrange'a--Sylwestra.
% Wyjaśnienie: Stopień jest mniejszy od $m$ ponieważ wyżej zdefiniowaliśmy $g$ dla wielomianu minimalnego ze stopniem $m$, a ponieważ L-S jest definiowany dla (m - 1) pochodnych, to musi być ostro mniejszy  

Otrzymujemy zatem następującą definicję funkcji od macierzy:
\begin{definition}[Funkcja od macierzy]
    Niech $f(\lambda)$ będzie funkcją zdefiniowaną na spektrum macierzy $A$, a $r(\lambda)$ --- odpowiednim wielomianem interpolacyjnym Lagrange'a--Sylwestra, wtedy:
    $$ f(A) := r(A) $$
\end{definition}

Udowodnimy ważną własność funkcji od macierzy:
\begin{theorem}\label{theorem:similar_matrix_functions}
    Jeśli macierze $A$ i $B$ są podobne, a macierz $T$ jest macierzą przejścia od $A$ do $B$:
    $$ B = T^{-1}AT $$
    wtedy macierze $f(A)$ oraz $f(B)$ również są podobne i macierz $T$ przekształca $f(A)$ w $f(B)$:
    $$f(B) = T^{-1}f(A)T$$
\end{theorem}
\begin{proof}
Najpierw pokażemy, że dwie macierze podobne mają ten sam wielomian minimalny:

Z $B = T^{-1}AT$ wynika:
$$
B^{k} = (T^{-1}AT)^{k} = T^{-1}AT \times T^{-1}AT \times \ldots \times T^{-1}AT = T^{-1}A^{k}T \quad (k \in \N)
$$
Wtedy dla dowolnego wielomianu:
$$
    g(\lambda) = a_0\lambda^{n} + a_1\lambda^{n - 1} + \ldots + a_{n - 1}\lambda + a_n 
$$
zachodzi:
$$
g(B) = a_0B^{n} + a_1B^{n - 1} + \ldots + a_{n - 1}B + a_nI
$$
Podstawiając $B^{k} = T^{-1}A^{k}T$:
$$
g(B) = a_0T^{-1}A^{n}T + a_1T^{-1}A^{n - 1}T + \ldots + a_{n - 1}T^{-1}AT + a_nI
$$
Z każdego składnika z lewa możemy wyciągnąć $T^{-1}$ oraz z prawa $T$:
$$g(B) = T^{-1}(a_0A^{n} + a_1A^{n-1} + \ldots + a_{n - 1}A + a_nI)T$$
Więc:
\begin{equation}\label{eq:local_equal_poly}
    g(B) = T^{-1}g(A)T
\end{equation}
Z tego wynika, jeśli $g(\lambda)$ jest wielomianem minimalnym dla $A$ to jest wielomianem minimalnym i dla $B$:
$$ g(A) = 0 \implies g(B) = T^{-1}g(A)T = 0$$
Odwrotnie: Z $B = T^{-1}AT$ mamy też $A = TBT^{-1}$, więc analogicznie do \eqref{eq:local_equal_poly}:
$$g(A) = Tg(B)T^{-1},$$ 
skąd $g(B) = 0 \implies (A) = 0$.

Udowodniliśmy, że dwie podobne macierze mają ten sam wielomian minimalny, a więc i te same wartości własne do rzędu $m_k - 1$. Z tego mamy, że funkcja $f(\lambda)$ przyjmuje te same wartości na spektrum maicerzy $A$ oraz na spketrum macierzy $B$ (Patrz Definicje~\ref{def:matrix_spectrum} \textit{Spektrum macierzy}). Więc istnieje wielomian interpolacyjny Lagrange'a--Sylwestra, taki że $f(A) = r(A)$ oraz $f(B) = r(B)$. Wtedy na mocy równości \eqref{eq:local_equal_poly} $g(B) = T^{-1}g(A)T$ mamy:
$$f(B) = T^{-1}f(A)T.$$
\end{proof}

Teraz udowodnimy kilka twierdzeń, które pozwolą nam rozszerzyć tożsamości znane dla funkcji jednej zmiennej na przypadek, gdy argumentem jest macierz.

\begin{lemma}\label{lemma:zero_spectrum_function}
Jeśli wielomian $g(\lambda)$ jest zdefiniowany na spektrum macierzy $A$ i wszystkie wartości na spektrum zerują się, tzn. $f(\Lambda_A) = 0$, wtedy wielomian $g(\lambda)$ jest podzielny przez wielomian minimalny $\phi(\lambda)$ macierzy $A$.
\end{lemma}
\begin{proof}
Z załozenia $g(\Lambda_A) = 0$, zapiszmy to jako w Definicji~\ref{def:function_on_spectrum} (\textit{Spektrum macierzy}):
$$ g(\lambda_k) = g'(\lambda_k) = \ \ldots \ = g^{(m_k - 1)}(\lambda_k) = 0, \quad (k = 1, \ldots, s) $$
Z tego wynika, że dla każdej wartości własnej $\lambda_k$ zachodzi:
$$ (\lambda - \lambda_k)^{m_k} \ | \ g(\lambda) $$
Ponieważ wartości własne są różne, czynniki $(\lambda - \lambda_k)^{m_k}$ są względnie pierwsze, więc:
$$ (\lambda - \lambda_1)^{m_1}\ldots(\lambda - \lambda_s)^{m_s} \ | \ g(\lambda) $$
Jak widzimy lewa strona jest właśnie wielomian minimalnym, więc
$$ \phi(\lambda) \ | \ g(\lambda) $$
\end{proof}

\begin{theorem}\label{theorem:base_matrix_properties}
Niech $G(u_1, u_2, \ldots, u_l)$ wielomian od $u_1, u_2, \ldots, u_l$ oraz $f_1(\lambda), f_2(\lambda), \ldots, f_l(\lambda)$ funkcji zdefiniowane na spektrum macierzy $A$. Zdefiniujmy następujący wielomian:
$$ g(\lambda) := G(f_1(\lambda), f_2(\lambda) \ldots, f_l(\lambda)).$$
Wtedy jeśli $g(\Lambda_A) = 0$, to $G(f_1(A), f_2(A), \ldots, f_l(A)) = 0$.
\end{theorem}
\begin{proof}
    Oznaczmy przez $r_1(\lambda), r_2(\lambda), \ldots, r_l(\lambda)$ wielomiany Lagrange'a--Sylwestra dla $f_1(\lambda), f_2(\lambda), \ldots, f_l(\lambda)$.

    Niech:
    $$ h(\lambda) := G(r_1(\lambda), r_2(\lambda) \ldots, r_l(\lambda)) $$
    Pokażemy, że jeśli $g(\Lambda_A) = 0$ to i $h(\Lambda_A) = 0$. Jak wiemy, z definicji wielomianu Lagrange'a--Sylwestra $f_j(\lambda_k)$ oraz $r_j(\lambda_k)$ mają te same wartości i pochodne do rzędu $m_k - 1$ na spektrum macierzy $A$, gdzie $j = (1, \ldots, l)$. Wielomian $G$ składa się z operacji dodawania oraz mnożenia (potęgowanie jest szczególnym przypadkiem mnożenia), na ogół każdy wielomian składa się z tych działań. Obie operacje zachowują zgodność pochodnych do rzędu $m_k-1$:
    pochodna sumy jest sumą pochodnych (liniowość), a pochodna iloczynu wyraża się wzorem Leibniza\cite[Rozdział III, \S 4, str. 205]{fichtenholz1999}:
    $$(ab)^{(n)}=\sum_{i=0}^{n}\binom{n}{i}a^{(i)}b^{(n-i)}.$$
    Z tego wynika, że $g$ i $h$ mają te same wartości i pochodne do rzędu $m_k-1$ w każdym $\lambda_k$, czyli $h(\Lambda_A) = 0 \implies g(\Lambda_A) = 0$.
    
    Z Lematu~\ref{lemma:zero_spectrum_function} wynika, że wielomian $h$ jest podzielne przez wielomian minimalny macierzy $A$, a więc jest wielomianem zerującym (Twierdzenie~\ref{theorem:annihilating_polynomial} \textit{O wielomianie zerującym}):
    $$ h(A) = G(r_1(A), r_2(A) \ldots, r_l(A)) = G(f_1(A), f_2(A) \ldots, f_l(A)) = 0$$
    % Note: Możemy podstawić $A$ ponieważ wielomian ma tylko operacji mnożenia i dodawania, co w przypadku podstawania macierzy jest homomorfizmem
\end{proof}

Korzystając z udowodnionego wyżej twierdzenia, możemy rozszerzyć tożsamości dotyczące funkcji zmiennej na analogiczne tożsamości dla funkcji od macierzy. Dalej przedstawiamy kilka przykładów niezbędnych w dalszych częściach pracy:

(\textit{1}) Niech $A$ macierz nieosobliwa, tzn. $\det{A} \ne 0$ oraz $f(\lambda) = \sqrt{\lambda}$. Niech $\sqrt{\lambda}$ będzie zdefiniowana na obszarze nie zawierającego zero oraz zawierającego wszystkie wartości własne macierzy $A$. Wtedy $\sqrt{\lambda}$ jest zdefiniowana na spekrum macierzy $A$, ponieważ z $\det{A} \ne 0$ wynika, że zero nie jest wartością własną oraz dla każdej wartości własnej $\lambda_k$ funkcja $\sqrt{\lambda}$ różniczkuję się do rzędu $m_k - 1$. Wtedy $f(A) = \sqrt{A}$ ma sens.

Niech $g(\lambda) = \lambda$ oraz
$$G(u_1, u_2) = u_1^2 - u_2,$$
wtedy
$$
    G(f(\lambda), g(\lambda)) = (\sqrt{\lambda})^2 - \lambda,
$$
Zauważmy, że funkcja $g(\lambda)$ jest zdefiniowana na spektrum macierzy $A$, więc na mocy Twierdzenia~\ref{theorem:base_matrix_properties} możemy powyższe równanie przyrównać do zera oraz podstawiając $G(f(A), g(A))$ mamy:
\begin{equation}\label{eq:sqrt_matrix}
    (\sqrt{A})^2 = A
\end{equation}

(\textit{2}) Niech $f(\lambda) = 1/\lambda$ oraz $A$ macierz nieosobliwa. Wtedy funkcja $f(\lambda)$ jest zdefinionowana na spektrum macierzy $A$. 

Niech
$$G(u_1, u_2) = u_1u_2 - 1,$$
wtedy
$$G(\lambda, f(\lambda)) = \lambda f(\lambda)  - 1$$
Przyrównując $G(\lambda, f(\lambda))$ do zera oraz na mocy Twierdzenia~\ref{theorem:base_matrix_properties} mamy:
$$ Af(A) - I = 0$$
$$ Af(A) = I$$
Ponieważ macierz $A$ jest nieosobliwa, to $A$ jest odwracalna, więc:
\begin{equation}\label{eq:invert_matrix_function}
    f(A) = A^{-1}
\end{equation}

Niech $r(\lambda)$ wielomian Lagrange'a--Sylwestra dla funkcji $f(\lambda) = 1/\lambda$, wtedy macierz $A^{-1}$ możemy przedstawić jako:
\begin{equation}\label{eq:invert_matrix_to_poly}
    A^{-1} = r(A)
\end{equation}
% Note: Ponieważ f(A) = r(A)
